% Figure: the principle of virtual work applied to a bell-crank lever.
% A rigid lever pivots at O; an input force P at arm a and an output load
% Q at arm b. A virtual rotation d-phi sweeps the two application points
% through virtual displacements a d-phi and b d-phi; setting the total
% virtual work to zero recovers the moment balance without solving for
% the pin reaction.
\begin{figure}[htbp]
\centering
\begin{tikzpicture}[>=latex, x=1.0cm, y=1.0cm]
  \coordinate (O) at (0,0);
  % the two arms in the reference configuration
  \coordinate (A) at (-2.6,0);   % input arm, length a
  \coordinate (B) at (0,-2.2);   % output arm, length b (bell crank, right angle)
  \draw[engblue, very thick] (A) -- (O) -- (B);
  % pivot
  \fill[engblack] (O) circle (1.8pt);
  \node[font=\scriptsize, anchor=south west] at (O) {$O$};
  \draw[engblack, thick] (O) -- (-0.28,0.42) -- (0.28,0.42) -- cycle;
  \draw[enggray] (-0.4,0.42) -- (0.4,0.42);
  % input force P at A, downward
  \draw[->, engred, very thick] (A) -- ($ (A) + (0,-1.0) $);
  \node[engred, font=\scriptsize, anchor=east] at ($ (A) + (-0.05,-0.5) $) {$P$};
  % output load Q at B, to the right
  \draw[->, engamber, very thick] (B) -- ($ (B) + (1.0,0) $);
  \node[engamber, font=\scriptsize, anchor=north] at ($ (B) + (0.5,-0.05) $) {$Q$};
  % arm labels
  \node[enggray, font=\scriptsize, anchor=south] at (-1.3,0.05) {$a$};
  \node[enggray, font=\scriptsize, anchor=west] at (0.05,-1.1) {$b$};
  % virtual rotation d-phi about O
  \draw[->, englime!70!black, thick] (0.55,0.0) arc (0:-32:0.55);
  \node[englime!55!black, font=\scriptsize] at (0.95,-0.35) {$\delta\phi$};
  % virtual displacement of A: perpendicular to OA, i.e. vertical here, length a d-phi
  \draw[->, englime!70!black, thick, dashed] (A) -- ($ (A) + (0,0.7) $);
  \node[englime!55!black, font=\scriptsize, anchor=south] at ($ (A) + (0,0.7) $) {$a\,\delta\phi$};
  % virtual displacement of B: perpendicular to OB, i.e. horizontal here, length b d-phi
  \draw[->, englime!70!black, thick, dashed] (B) -- ($ (B) + (0.7,0) $);
  \node[englime!55!black, font=\scriptsize, anchor=west] at ($ (B) + (0.7,0) $) {$b\,\delta\phi$};
\end{tikzpicture}
\caption[Virtual work on a bell-crank lever]{The principle of virtual
work read off a bell-crank lever pivoted at $O$. An input force $P$ acts
at arm $a$, an output load $Q$ at arm $b$. Give the rigid lever a virtual
rotation $\delta\phi$ about the pivot. Each application point sweeps an
arc perpendicular to its arm, of length $a\,\delta\phi$ and $b\,\delta\phi$
respectively, and the pin reaction at $O$ does no work because $O$ does
not move. Setting the total virtual work to zero,
$P\,a\,\delta\phi - Q\,b\,\delta\phi = 0$, gives $Pa = Qb$, the moment
balance, without ever computing the pin reaction. The virtual-work route
delivers the input-output relation of any single-degree-of-freedom
mechanism in one equation, sidestepping the connection forces that the
free-body method would have to carry.}
\label{fig:vol03ch02:virtual-work-lever}
\end{figure}
